\section{微分公式を逆向きに読む}\label{節:基本公式の逆}

$F'(x)=f(x)$ ならば，$F$ は $f$ の原始関数である。この関係を使うとき，最初から正しい係数まで思い付く必要はない。まず形の近い候補を微分し，元の関数との差を直すのでよい。たとえば
\[
 \frac{\dd}{\dd x}x^5=5x^4
\]
だから，$x^4$ の原始関数を得るには，$x^5$ を $5$ で割る。
\[
 \int x^4\dd x=\frac{x^5}{5}+C.
\]
同じ理由で，$r\ne-1$ のとき
\[
 \int x^r\dd x=\frac{x^{r+1}}{r+1}+C
\]
となる。指数を一つ増やし，増やした指数で割るのは，微分で指数が一つ減り，その指数が係数として出ることを元に戻しているのである。

ここで $r=-1$ だけは別に考える必要がある。上の式へ $r=-1$ を代入すると分母が $0$ になり，意味を持たない。また，$x^0=1$ を微分しても $x^{-1}$ は現れない。代わりに，対数の微分
\[
 \frac{\dd}{\dd x}\log x=\frac1x\qquad(x>0)
\]
を使う。$x<0$ では $\log(-x)$ を微分すれば
\[
 \frac{\dd}{\dd x}\log(-x)
 =\frac1{-x}\,(-1)=\frac1x
\]
となるから，二つの形は $\log|x|$ とまとめられる。

\begin{theorem}{基本的な不定積分}
次の公式は，右辺が微分でき，左辺の関数が定義される区間で成り立つ。
\[
\begin{array}{c|c}
 f(x)&\int f(x)\dd x\\ \hline
 x^r\quad(r\ne-1)&\frac{x^{r+1}}{r+1}+C\\[5pt]
 \frac1x&\log|x|+C\\[5pt]
 \e^x&\e^x+C\\[3pt]
 a^x\quad(a>0,\ a\ne1)&\frac{a^x}{\log a}+C\\[5pt]
 \sin x&-\cos x+C\\[3pt]
 \cos x&\sin x+C\\[3pt]
 \frac1{\cos^2x}&\tan x+C
\end{array}
\]
実数 $r$ に対するべきの公式は，まず $x>0$ で考える。整数のべきなど，負の $x$ でも定義される場合は，その定義域内の区間にも適用できる。
\end{theorem}

この表は，右の式を微分して確かめられる。たとえば $a^x$ の微分では $\log a$ が掛かるので，積分では $\log a$ で割る。$a=1$ を除くのは，この割り算ができないためであり，$1^x=1$ の積分そのものは $x+C$ である。

不定積分の $C$ は，考えている一つの区間上での任意定数である。$\frac1x$ の定義域は $0$ の左右に分かれる。左側で加える定数と右側で加える定数が，同じである必要はない。図\ref{図:対数の区間}では，この区間の分かれ方を数直線に示した。原始関数を一つ選ぶ場合と，定義域全体のすべての原始関数を述べる場合とを区別しよう。

\begin{figure}[htbp]
 \centering
 \begin{tikzpicture}[x=1cm,y=1cm]
  \draw[-{Stealth}] (-4,0)--(4,0) node[right] {$x$};
  \draw[very thick,blue] (-3.7,0)--(-0.12,0);
  \draw[very thick,blue] (0.12,0)--(3.7,0);
  \filldraw[fill=white] (0,0) circle (2pt);
  \node[below] at (0,-0.08) {$0$};
  \node[above] at (-2,0.18) {$\log(-x)+C_-$};
  \node[above] at (2,0.18) {$\log x+C_+$};
  \node[below] at (-2,-0.35) {$x<0$};
  \node[below] at (2,-0.35) {$x>0$};
 \end{tikzpicture}
 \caption{$0$ を越えて一つの区間にはできない。定数は左右で独立に選べる。}
 \label{図:対数の区間}
\end{figure}

\begin{example}{根号と分母をべきの形へ直す}
$x>0$ において，$\int\left(2\sqrt{x}+\frac3{x^2}\right)\dd x$ を求めよ。
\end{example}

根号のままでは公式と見比べにくいので，$\sqrt{x}=x^{\frac12}$，$\frac1{x^2}=x^{-2}$ と書く。二つの項は別々に積分できるから，
\[
\begin{aligned}
 \int\left(2\sqrt{x}+\frac3{x^2}\right)\dd x
 &=2\int x^{\frac12}\dd x+3\int x^{-2}\dd x\\
 &=2\cdot\frac23x^{\frac32}
   +3\cdot\frac{x^{-1}}{-1}+C\\
 &=\frac43x^{\frac32}-\frac3x+C.
\end{aligned}
\]
分母にある式をいつも対数で積分するわけではない。$\frac1x$ は指数が $-1$ なので対数になり，$\frac1{x^2}$ は指数が $-2$ なので，べきの公式を使う。答えを微分すると
\[
 \frac43\cdot\frac32x^{\frac12}
 +3x^{-2}=2\sqrt{x}+\frac3{x^2}
\]
となり，係数と符号の両方を確認できる。

和を項ごとに積分できるのは，微分が和と定数倍を保つためである。一方，積について
\[
 \int f(x)g(x)\dd x
 =\left(\int f(x)\dd x\right)
  \left(\int g(x)\dd x\right)
\]
という規則はない。右辺を微分すると積の微分の二項が現れ，ふつうは $f(x)g(x)$ にはならない。積があるときは，まず展開できるか，恒等式で和にできるかを考える。その後に学ぶ置換積分と部分積分も，この問題への別の答えである。

